% !BIB program = biber
\documentclass[stu]{apa7}

\raggedbottom

\usepackage[american]{babel}
\usepackage{lipsum}
\usepackage{csquotes}
\usepackage{booktabs}
\usepackage[style=apa,backend=biber]{biblatex}

\addbibresource{ppb.bib}
% \bibliographystyle{unsrt} % Or use a specific apa-numeric style
% \bibliography{ppb.bib}

%% --- Proposal Metadata ---
\title{Physics Informed Hyperspectral Image Reconstruction from RGB}
\shorttitle{Project Proposal} % Required for APA headers

\authorsnames{Aditya Bhargava}
\authorsaffiliations{University of Georgia}
\course{CSCI 8820: Computer Vision and Pattern Recognition}
\professor{Dr. Suchendra Bhandarkar}
\duedate{March 18, 2026}

%% --- Abstract (Proposal Executive Summary) ---
\abstract{This proposal outlines the development of a system, aiming to solve HSI reconstruction. We discuss the motivation, the proposed technical approach using computer vision techniques, and the expected deliverables.}

\keywords{Computer Vision, HSI, RGB, Physics Informed}

\begin{document}
\maketitle



\section{Introduction and Problem Statement}

% Describe the "Why". What is the gap in current computer vision research or applications?
Hyperspectral imaging (HSI) is an advanced optical sensing technique that assembles spectroscopy and digital photography into a single system  \parencite{hong2026hyperspectralimaging}. The imaging collects and processes scene information by dividing the whole spectrum into tens of hundreds of bands, typically covering wavelengths from 380 nm to 2500 nm \parencite{Rybicki:847173}. The technique has been widely adopted in many fields, eg., medical diagnosis, healthcare, remote sensing and identifying foreign objects in food products \parencite{ghimpeteanu2025hyperspectralimagingidentifyingforeign}. The hyperspectral images possess a three-dimensional structure, comprising two dimensions of spatial information and one dimension of spectral information. However, conventional spectrometers often operate in 1D line or 2D plane scanning, which is time-consuming and not suitable for dynamic scenes. Also, the devices for acquring the HSIs are typically more complicated and expensive than common cameras. 
% In order to obtain the spectral information more effectively by extracting the required spectral information from the RGB images. 
The need to reconstruct the spectral properties of scene reflectance (and illumination) from a three-channel device or standard color space was recognized as early as 1980s \parencite{Hall1988IlluminationAC,Jaaskelainen:90, Maloney:86, Maloney:86, Parkkinen:89}. But this is an extremely ill-posed problem since a large amount of information is lost during the process of spectral integration when RGB sensors capture the light. Spectral reconstruction from RGB images often suffer from a \textit{metameric} dilemma, where distinct spectral distribution map to nearly identical RGB values, making them indistinguishable. Recent approaches involve deep learning methods to solve this by learning the mapping from RGB to HSI. While these models achieve reasonable quantitative metrics, they operate as black boxes that ignore the underlying mechanism behind the imaging process. To overcome these issues, the goal is to incorporate physics based formulations into the architecture itself. By explicitly embedding the camera's Spectral Response Function (SRF) as a mathematical forward-model constraint, the network is forced to output hyperspectral reconstructions that strictly adhere to the physical laws of image formation, achieving both higher accuracy and reliable material distinction.

\section{Literature Review}
A RGB image is usally generated in three wide spectral channels with Bayer Filters, which is essentially a color filter array (CFA) placed over digital image sensor to capture color information. While, an HSI is obtained by the interaction between the spectral reflectance and illumination. The relation between the RGB and HSI images can be represented by

\begin{equation}
    I_c(x, y) = \int_{w_1}^{w_2} R(x, y, w) L(w) S_c(w) \, dw
\end{equation}
where $I_c$ represents each channel of RGB image, $S_c()$ denotes the spectral response function and $R(x,y,z)$ is the spectral reflectance of the objects and $L(w)$ is the illumination spectrum. HSI is typically defined as
\begin{equation}
    I_c(x, y) = \int_{w_1}^{w_2} H(x, y, w) S_c(w) \, dw
\end{equation}
where, $H(x,y,z)$ represents the Hyperspectral image. The challenge in spectral reconstruction is to invert the above forward image formation model in optimzation framework to find the $H(x, y, z)$. A wide variety of methods have been used to achieve the reconstruction.

\subsection{Prior Methods}
These methods rely on mathematical assumptions around which the algorithms are designed. Broadly there are three categories of Prior-Based methods namely dictionary learning, manifold learning and gaussian process. \textcite{Arad2016SparseRO} introduced a method to reconstruct HSI directly from RGB leveraging hyperspectral priors in order to create a sparse dictionary of hyperspectral signatures and their corresponding RGB projections. The method works by collecting the hyperspectral prior, preferably from a set of domain specific scenes. This prior is then reduced computationally to an overcomplete dictionary $D_h$ of hyperspectral signatures. The dictionary is projected to the sensor space via the receptor spectral absorbance functions $R$, establishing a corresponding RGB dictionary, $D_{rgb}$. 
This correspondance between each RGB vector and its hyperspectral signature originator $h_i$ is maintained for the later mapping from RGB to hyperspectral signatures. For a given image, a sparse weight vector $w$ is calculated, that reconstructs a given RGB pixel query from the RGB dictionary. Because both dictionaries share the exact same sparse coefficients, applying these weights back to the hyperspectral dictionary yields the estimated spectrum signature $h_q$. While, in manifold learning approach is aimed at finding a low-dimensional manifold 

% this formulation inherently guarantees that the reconstructed hyperspectral signal is physically consistent with the original RGB measurement (i.e., $c_q = R \cdot h_q$), though the overall reconstruction quality remains strictly bounded by the representational capacity of the initial hyperspectral prior.

\subsection{Deep-Learning Methods}
\section{Proposed Methodology}
% How are you going to build it?
\subsection{Data Acquisition and Preprocessing}
The dataset that will be used is called Arad 1K as it also contains the camera sf matrix

\subsection{Proposed Algorithm/Model Architecture}


\subsection{Evaluation Metrics}
Accuracy, time taken 

\section{Project Plan and Timeline}
The project will be executed in four phases as detailed in Table~\ref{tab:Timeline}.

\begin{table}[h]
  \caption{Project Milestone Timeline}
  \label{tab:Timeline}
  \begin{tabular}{@{}lll@{}} \toprule
  Phase & Task Description & Completion Date \\ \midrule
  1     & Data Collection \& Cleaning & March 25      \\
  2     & Initial Model Training      & April 10      \\
  3     & Testing and Optimization    & April 25      \\
  4     & Final Report \& Submission  & May 05        \\ \bottomrule
  \end{tabular}
\end{table}

\section{Expected Outcomes and Impact}


\printbibliography

\appendix

\section{Resource Requirements}
\label{app:resources}
% List hardware (GPUs), datasets, or software libraries needed.
\begin{figure}[h]
    \caption{Proposed System Architecture Flowchart.}
    \figurenote{This diagram represents the data pipeline from input sensor to classification.}
    \label{fig:Arch}
\end{figure}

\end{document}